% NOTE (2026-09-25): superseded by ../draft.tex, which was revised after a 46-item review; this file is kept as a source only.
% =====================================================================================
% sections/discussion.tex -- Discussion, next steps and Limitations of the short draft of "Degrees of Separation"
% (author: Sora Yng). Written 2026-09-25 for \input; no preamble. Terminology follows theory_core.tex.
% Every number is copied from paper/main.tex (Results, Table "ladder", Limitations) or theory_core.tex, which copies
% main.tex, si.tex and THEORY_TESTS_RESULT.md; none is new. Citation keys: ../refs_merged.bib, plus anelli2020returns
% from refs_short.bib (verified against Crossref, DOI 10.1093/jeea/jvz070).
% =====================================================================================
\providecommand{\ci}[2]{$[#1,\allowbreak\,#2]$}

\section{Discussion and next steps}\label{sec:discussion}

\paragraph{For students choosing a programme.} Where the academic and the graduate labour markets order the same
programmes alike, they agree at the university's name. Most of that agreement is shared with the ordering of intakes
that admission produces: SAT and admission rate take mean US coupling from $+0.43$ to $+0.14$. Below the name the
two markets agree little, clearly only in computer science. Within a university, a department one standard deviation
higher in $F$ goes with about 0.8\% higher four-year pay ($+0.0084$ log points, \ci{+0.0000}{+0.0164}). For a
student choosing among university $\times$ field programmes, a department's academic reputation thus adds little to
what the name and the intake already say about graduates' pay. It is a weak guide rather than an empty one, because
error in $F$ dilutes the department part, which may reach about 3\% under the lower-bound reliability. How strongly
the name shows in pay depends on the field, and coupling is weak where pay follows national scales. None of this
says what a programme would add to a given student's pay.

\paragraph{For rankings.} The same bound applies to subject rankings built on academic standing. In our data the
field rank $F$ plays that part, and it predicts graduates' pay better than $G$ in only 15 of 57 fields.
\citet{britton2022how} report a similar pattern for UK league-table ratings net of intake, which suggests that it is
not specific to a hiring-based measure. Because rankings move applications \citep{bowman2009getting,
luca2013salience}, a subject ranking that mostly restates the name may steer applicants between names while
appearing to compare programmes. Rankings by graduates' pay order intakes as well as programmes, because a
programme's median includes whom it admits \citep{dale2002estimating}. Under national pay scales they reveal little
about employers' valuation.

\paragraph{From agreement to returns.} Agreement between orderings is not a return, since the name term in coupling
can be a pure common cause (Section~\ref{sec:thy}). Much of the premium of selective colleges reflects whom they
admit \citep{dale2002estimating, dale2014estimating}, and their causal effects concentrate in the upper tail that
medians miss \citep{zimmerman2019elite, chetty2023diversifying}. The open question is whether an elite-university
return belongs to the university's name or to the programme. Admission cutoffs could answer it, because applicants
just either side of a programme's cutoff differ in admission but barely in the score that ranks them. The design
therefore holds intake ability fixed at the margin, where our aggregates can only partial out noisy proxies of it.
Cutoff designs have estimated returns to a flagship state university \citep{hoekstra2009effect}, to a highly
selective university \citep{anelli2020returns}, and to fields relative to applicants' next-best choice
\citep{kirkeboen2016field}. \citet{anelli2020returns} also finds that part of the elite enrolment premium runs
through the different fields its students choose. At the university $\times$ field level the same logic splits a
return in two. Applicants whose next-best option is the same field at another university identify a contrast between
universities within the field. Applicants whose next-best option is another field at the same university identify a
contrast between fields that holds the name fixed. Across many contrasts of the first kind, the part of the return
that follows the universities' gap in $G$, field by field, belongs to the name. The part that follows the
departments' gap in $F$ beyond it belongs to the programme. This is the causal counterpart of our comparison of $F$
with $G$, local to applicants at the cutoff.

\paragraph{What the design needs.} For each applicant, the design needs the admission score, the cutoff and offer of
each programme applied to, and the ranked list of programmes, which identifies the next-best option. The ranked
lists would also measure students' ordering of programmes from their own choices \citep{avery2013revealed}, at the
university $\times$ field level where our selectivity measures stop at the university. It also needs enrolment,
field switches, completion, and earnings from tax or social-security records, with their whole distribution and
where graduates work. It needs these earnings for a decade or more: first-year pay cannot be read as employer
reputation here, and within fixed cohorts US coupling rises by $+0.031$ per year \ci{+0.021}{+0.041}. Cutoffs must
bind in enough programmes per university and per field to compare them. Finally, the same departments need an
academic reputation, which outside the United States would have to be built from faculty-hiring records, as we built
the UK hierarchy \citep{li2026orcid}.

\subsection{Limitations}\label{sec:limits}

The analysis is descriptive and uses only public aggregates. It does not separate value added, selection and
employer screening, does not measure beliefs about prestige, and cannot say why agreement below the name is weak.
Programme medians miss the tails, and Scorecard medians cover only Title~IV-aided completers who work and are not
enrolled. Selectivity is measured for the university, not the programme, and whatever the proxies miss stays in
every partial coupling. $F$ is a net-placement position that is partly doctoral production, a share the audit bounds
only from above. Its reliability can only be bracketed between 0.81 and 0.96. Only one check nets pay of destination
location, and only between census divisions. About a fifth of pay coupling runs between employer sectors, which
reflects employers' valuation only if employers ration access to sectors. The UK hierarchy rests on self-reported
ORCID careers of uncertain validity. The theory was written after most results were known; it organises them, and
they are not evidence for it. Its nine tests were specified before the analysis; not publicly registered. They share
data with results already seen, and two of them are inconclusive for lack of power.
